% Storyline 2, paragraph 3: from information synchronization to individual integration of evidence.
Most proposed remedies keep punishers in good standing by aligning what observers know, because cooperation is stable only when individual opinions are sufficiently correlated\cite{murase2024synchronization}. Public institutions broadcast a single verdict\cite{radzvilavicius2021}, gossip spreads opinions between peers\cite{sommerfeld2007,kawakatsu2024}, and other rules synchronize private opinions directly\cite{murase2024synchronization} or evaluate actions from the donor's perspective\cite{radzvilavicius2019}. Each route rests on conditions that are often unmet. Institutions need monitoring capacity and the adherence of those they serve\cite{radzvilavicius2021}, gossip must be frequent and accurate enough to produce consensus\cite{kawakatsu2024,seki2016}, and empathetic evaluation requires knowing how the donor sees the recipient\cite{radzvilavicius2019}. These conditions are strained where empathy, communication and institutional trust are compromised, as in polarized societies\cite{iyengar2019,michelmata2024}. Recent models therefore seek agreement from private information alone, by letting each observer judge from several fresh observations of the same individual\cite{michelmata2024} or assess recipients as well as donors\cite{tham2026}. Both approaches enlarge the evidence behind each judgement. Whether agreement can also arise from how individuals integrate the evidence they already receive, one observation at a time, remains open. In principle it can, because independent observers agree whenever each reaches a reliable verdict on its own. A rule that makes verdicts reliable by discounting contrary evidence, however, may also discount genuine defection.
